\documentclass[11pt]{exam}
\usepackage{enumerate}
\usepackage[utf8]{inputenc}
\usepackage{hyperref}
\usepackage{alltt}
\usepackage{xspace}
\usepackage{longtable}
\usepackage{amsfonts}
\usepackage{latexsym}
\usepackage{amsmath}
\usepackage{amssymb}
%\usepackage{xcolor}
\usepackage{enumerate}
\usepackage{graphicx}
%\usepackage{moreverb}
%\usepackage{ifthen}
\usepackage{url}
\usepackage{tikz} 
\usepackage{bussproofs}

%\usetikzlibrary{snakes}
\usetikzlibrary{decorations}
\usetikzlibrary{shapes,shapes.geometric,arrows,fit,calc,positioning,automata,decorations.pathmorphing}
\tikzset{state/.style={draw,ellipse}}

\usepackage{algorithmicx}
\usepackage[noend]{algpseudocode}
\usepackage[Algoritmo]{algorithm}


\input sli_aux
\newboolean{portuguese}
\setboolean{portuguese}{true}
\setboolean{portuguese}{false}

\newboolean{normal}
\setboolean{normal}{true}

\newcommand{\pten}[2]{\ifthenelse{\boolean{portuguese}}{#1}{#2}}
\pten{\usepackage[portuges]{babel}}{\usepackage[british]{babel}}

%\pten{\uselanguage{portuguese}}{\uselanguage{british}}

\setcounter{aula}{3}

\begin{document}
	\begin{center} 
{\bf {\large \pten{Programação Concorrente}{Concurrent Programming} - Exercícios \theaula} \\
\medskip

CCS  }
\end{center}
%\pten{Nota: perguntas com (**) não facultativas}{Questions with (**) are optional}.
\begin{questions}
\question \pten{Indica quais das seguintes expressões são expressões do CCS sintaticamente correctas}{Which of the following expressions of CCS are correct?}.
\begin{parts}
\part $(a.0+\overline{a}.A)\backslash \{a,b\}$
\part 	 $(a.0+\overline{a}.A)\backslash \{a,\tau\}$
%\part  (**) $a.B+ [a/b]$
\part $\tau.\tau . B+ 0$
%\part (**) $(a.B+b.B)[a/b,b/a]$
%\part (**) $(a.B+b.B)[a/\tau,b/a]$
\part $(a.b.B+\overline{a}.0) | B$
\part $(a.b.B+\overline{a}.0) . B$
\end{parts}
\question \pten{Sendo}{Let} $A:= b.a.B$, \pten{usando as regras de inferência mostra a existência das seguintes transições}{using  the inference rules show that the following transitions exist}:
\begin{itemize}
\item  $\trans{(A|\overline{b}.0)\backslash\{b\}}{\tau}{ (a.B|0)\backslash \{b\}}$	
%\item (**)$\trans{(A|\overline{b}.a.B)+(\overline{b}.A)[a/b]}{\overline{b}}{ (A|a.B)}$	
%\item (**)$\trans{(A|\overline{b}.a.B)+(\overline{b}.A)[a/b]}{\overline{a}}{ (A[a/b]}$	
\end{itemize}
	
\question \pten{Resolver os exercícios}{Solve the}
$CCS$ \pten{em}{ exercises in} \href{https://pseuco.com/#/exercises}{PseuCo.com}

\question
\pten{Considera as seguintes definições uma investigador numa universidade que toma café se publica artigos científicos}{Consider the following definitions of a researcher that takes coffee and publishes articles}.

\begin{eqnarray*}
	CM&:=& coin?.coffee!.CM\\
CS&:=&pub!.coin!.coffee?.CS\\
Uni&:=& (CM|CS)\backslash\{coin,coffee\}
\end{eqnarray*} 
\pten{Usa as regras de inferência do CCS para obter um fragmento atíngível  de}{Use the CSS inference rules to obtain the reachable fragment of} $\bras{Uni}_\Gamma$. \pten{Testa no}{Test in} pseuco.com.
\pten{Comparar com}{Compare with } $Spec:=pub!.Spec$.

\question \pten{Sendo}{Let} $A:=(a.A)\backslash\{b\}$ \pten{mostra que}{show that} $\bras{A}_\Gamma$ \pten{é infinito}{infinite} \pten{(mesmo o fragmento atingível)}{(even the reachable fragment)}.
\question 
\pten{Considera as seguintes definições que pretendem resolver o problema da exclusão mútua com um semáforo (unário)}{The  following definitions try to solve the mutual exclusion with a semaphore}.
\begin{parts}
\part \begin{eqnarray*}
	Mutex_1&:=& (User|Sem)\backslash \{p,v\}\\
User&:=&\overline{p}.enter.exit.\overline{v}.User\\
Sem&:=& p.v.Sem
\end{eqnarray*} 
\pten{Usa as regras de inferência do CCS para obter um fragmento atíngível  de}{Use the CCS inference rules to obtain the reachable fragment of} $\bras{Mutex_1}_\Gamma$. \pten{Testa no}{Test in} pseuco.com.
\part \pten{Seja}{Let}

\begin{eqnarray*}
	Mutex_2&:=& ((User|Sem)|User)\backslash \{p,v\}
\end{eqnarray*} 
\pten{Usa as regras de inferência do CCS para obter um fragmento atíngível  de}{Use the CCS inference rules to obtain the reachable fragment of} $\bras{Mutex_2}_\Gamma$. \pten{Testa no}{Test in} pseuco.com.
\pten{Havia alteração se}{There will be changes if one uses} $User := \overline{p}.enter.\overline{v}.exit.User$ ?
\part \pten{Seja}{Let}
\begin{eqnarray*}
	FMutex&:=& ((User|Sem)|FUser)\backslash \{p,v\}\\
	FUser &:=& \overline{p}.enter.(exit.\overline{v}.FUser +exit.\overline{v}.0)
\end{eqnarray*} 

\pten{Usa as regras de inferência do CCS para obter um fragmento atíngível  de}{Use the CCS inference rules to obtain the reachable fragment of} $\bras{FMutex}_\Gamma$. \pten{Testa no}{Test in} pseuco.com.
\pten{Achas que}{Do you think that} $Mutex_2$ \pten{e}{and} $FMutex$ \pten{têm o mesmo comportamento}{has the same behaviour}?
\end{parts}

\question 
Modela uma  passagem de nível com 3 componentes: o comboio, a cancela e um controlador (como sugerido na figura). Considera que:

\begin{itemize}
\item As acções do comboio  caracterizam a sua localização em relação à passagem: aproxima-se ($approach$), entra ($in$) e sai ($exit$).
\item A cancela pode estar baixada  ($down$) ou levantada ($up$).
  \item O comboio e a cancela não comunicam directamente, apenas o controlador comunica com eles. O controlador deve garantir que quando o  comboio entra a cancela está em baixo.
  \end{itemize}
  \begin{center}
%\includegraphics[width=200pt]{train.png}
[TRAIN]
\end{center}\begin{parts}
  \part Descreve em CCS  cada componente: comboio, cancela e controlador.
\part Escreve o processo do sistema que corresponda à execução paralela das várias componentes, indicando as acções em que devem sincronizar. 
  \part 
  \pten{Usa as regras de inferência do CCS (indicando em cada passo a regra usada) para obter um LTS para esse processo}{Use the CSS inference rules (mentioning the name of the rule used in each step) to obtain the reachable fragment of that process}.  \end{parts}
\question \pten{Considera os seguintes conjuntos de equações  que definem dois contador}{Consider}:
\begin{eqnarray*}
C_0&:= &inc. C_1\\	
 C_n&:=& inc.C_{n+1}+dec.C_{n-1}, \text{ para } n\geq 1\\
\end{eqnarray*}
e $SC:= inc.(SC|dec.0)$.
\begin{parts}
\part  \pten{Determinina o fragmento acessível  (e inicial) do LTS}{Determine the inicial fragment}  $\bras{C_0}_\Gamma$. \pten{Argumenta que o LTS não pode ser finito por estados}{Argue that the LTS cannot be state-finite}.
\part \pten{Determinina o fragmento acessível do LTS (e inicial) $\bras{SC}_\Gamma$}%
				   {Determine the accessible fragment of LTS (and initial) $\bras{SC}_\Gamma$}
\part \pten{Determinita se os dois LTS são isomorfos. Haverá alguma semelhança entre eles?}%
					 {Determine if the two LTS are isomorphic. Is there any similarity between them?}
\end{parts}
\end{questions}
\end{document}